\section{Extending the Tool}
\label{sec:extending}

Rules that keep extensions safe: every input is a configuration parameter, every fixed number is in
\texttt{constants.py}, quantities are SI, every function has a docstring. After any change, run
\texttt{python tests/verify\_reference.py} (Section~\ref{sec:testing}).

\subsection{Adding a profile family}
\label{sec:newfamily}

A new member of an existing family needs no code (\key{nose.custom\_profiles}, Section~\ref{sec:structured}). A new
\emph{family} is one function, which can live in a plugin file (\texttt{plugins/secant\_ogive.py}):
\begin{lstlisting}[language=Python]
import numpy as np
from nosecone.geometry import register_family

@register_family("secant_ogive", params={"rho_factor": 1.5}, dimensional=True,
                 doc="Secant ogive: arc of radius rho_factor x the tangent-ogive radius")
def secant_ogive(x, L, R, rho_factor):
    rho = rho_factor * (R**2 + L**2) / (2 * R)
    alpha = np.arccos(np.sqrt(L**2 + R**2) / (2 * rho)) - np.arctan(R / L)
    x = np.clip(x, 0.0, L)
    return np.sqrt(np.maximum(rho**2 - (rho*np.cos(alpha) - x)**2, 0.0)) - rho*np.sin(alpha)
\end{lstlisting}
The function takes NumPy arrays and returns $y(0) = 0$, $y(L) = R$, monotonic (with the default
\texttt{dimensional=False} it takes $x/L$ and returns $y/R$). Load it with \key{project.plugins}, declare it in
\key{nose.custom\_profiles} and set \key{aerodynamics.compute\_missing: true} for its wave drag
(\texttt{configs/example\_plugin.yaml}); everything else follows automatically.

\subsection{Adding a plot}

A plot is a function of a \texttt{PlotContext} (\texttt{ctx.cfg}, \texttt{ctx.study}, \texttt{ctx.trade},
\texttt{ctx.traj}) registered in \texttt{PLOTS}:
\begin{lstlisting}[language=Python]
# nosecone/plotting/figures.py
def fig_mass(ctx: PlotContext) -> None:
    fig, ax = plt.subplots(figsize=(fs.WIDTH, 0.35 * fs.WIDTH))
    for n in ctx.order():
        c, ls, mk = ctx.sty(n)
        d = ctx.trade[ctx.trade["shape"] == n].sort_values("fineness")
        ax.plot(d.fineness, d.mass_g, color=c, ls=ls, marker=mk, label=ctx.label(n))
    ctx.save(fig, "fig_mass")

PLOTS["mass"] = fig_mass          # "mass" is now valid in output.plots
\end{lstlisting}

\subsection{Adding a parameter or changing a model}

\begin{itemize}[leftmargin=*,itemsep=1pt]
\item \textbf{A parameter}: one line in the dataclass of its section (\texttt{nosecone/config.py}), e.g.\
\lstinline[language=Python]|coupler_mass: float = param(0.120, "kg", "Bonded coupler mass", NONNEG)|; it is then
validated, listed by \texttt{describe-config} and available as \texttt{cfg.structure.coupler\_mass}.
\item \textbf{A coefficient} (friction, base drag, heating): \texttt{constants.py}.
\item \textbf{The drag build-up}: \texttt{flight.DragModel.components()}; a new component is added to the total,
scaled by $k$ and shown in \texttt{breakdown}.
\item \textbf{The transonic wave drag}: \texttt{aerodynamics.WaveDrag.\_\_call\_\_}; with RANS data, add the points
to the database and lower \key{aerodynamics.transonic\_match}.
\end{itemize}

\subsection{Checking a change}
\label{sec:testing}

\texttt{python tests/verify\_reference.py} reruns the Weisshorn study and compares every number, table and CAD file
with the reference outputs (\texttt{-{}-figures}, \texttt{-{}-cfd} add the figures and three CFD cases). A change not
meant to alter the results must pass unchanged; if it is meant to, update \texttt{tests/reference/} only after the new
results have been reviewed.
